% Generated by tools/md2tex.py from drafts/paper_en.md. Edit the draft, then re-run the script.
\section{Introduction}
\label{sec:introduction}

Subject-to-video (S2V) generation renders a user-provided subject, given as a reference image, into a video described by a text prompt. Recent systems \cite{liu2025phantom, jiang2025vace, fei2025skyreelsa2, hu2025hunyuancustom, deng2025magref, xue2025standin} obtain this ability by fine-tuning a pretrained text-to-video backbone \cite{wan2025, kong2024hunyuanvideo} on paired reference--video data.

These pairs are the bottleneck. They are scarce, because every clip needs a clean, subject-centric reference whose identity is verified across the clip. More importantly, they are biased. The reference is usually cropped from a frame of the target video, so it shares pose, viewpoint, lighting, and often background with the target. Training on such pairs rewards a trivial solution: paste the reference into the frame. Recent benchmarks \cite{yuan2025opens2v, muss2025} document the resulting \emph{copy-paste} behavior: the subject stays frozen in the reference pose, ignores prompted actions, and carries the reference lighting into the new scene.

Subject-driven image data suffer much less from either problem. Such datasets are orders of magnitude larger \cite{xiao2024omnigen, tan2024ominicontrol, wu2025uno}, and their reference--target pairs are constructed to differ in pose, scene, and viewpoint, so a model trained on them must re-render the subject rather than copy it. Our central proposition is therefore: \textbf{learn appearance from images, and let the video model contribute only motion.}

The obstacle is that images carry no motion. Supervising a video model with an image repeated into a static clip teaches it that the correct output is still, and the motion prior of the backbone is lost. Image-only training is useful only if static supervision can shape appearance without teaching the model to be static, and if appearance and motion remain separable afterwards.

\method{} meets both conditions on a frozen video diffusion transformer with two mechanisms (Fig.~\ref{fig:overview}).

\begin{itemize}
  \item \textbf{A motion carrier with a static regime.} A Stage-I low-rank adapter is trained with signed targets, static clips at $M{=}{+}1$ and real videos at $M{=}{-}1$, so that a runtime scale $M$ becomes a continuous motion control. During image training the carrier runs at $M{=}1$ and already produces the still video that the target demands. The appearance branch is therefore never asked to suppress motion, and setting $M{\le}0$ at inference restores it.
  \item \textbf{A paired-difference reference readout.} For every input, the appearance branch is evaluated once with the real reference and once with a learned null reference, and only the difference of the two predictions is added to the carrier. Everything that does not depend on the reference cancels, the carrier is never touched by reference-dependent computation, and the reference-free branch needed for guidance is learned by construction instead of by dropout.
\end{itemize}

As a result, motion, text adherence, and reference adherence become three independent controls at inference time.

Our contributions are as follows.

\begin{itemize}
  \item We trace copy-paste to how paired video data are constructed, and show that subject-to-video generation can be trained from image pairs alone, with the pretrained video model serving purely as a source of motion.
  \item We introduce a signed motion carrier whose static regime absorbs static image supervision, so that image training shapes appearance without erasing motion.
  \item We introduce a paired real-versus-null difference readout that structurally separates the reference signal from the motion computation and yields independent motion, text, and reference controls.
  \item We propose a copy-paste metric that is simple to compute and validated on synthetic paste and natural re-render anchors, and report it together with a same-subject gate, so that low copying is not mistaken for failing to render the subject.
\end{itemize}
